\chapter{Interface utilisateur}

\section{Vue d'ensemble}

L'interface principale est organisée autour de quatre zones :

\begin{itemize}
  \item l'arbre du modèle, à gauche ;
  \item la vue graphique principale, au centre ;
  \item le panneau de propriétés, à droite ;
  \item les onglets d'informations, de console et de résultats, en bas.
\end{itemize}

\captureplaceholder{Vue générale avec arbre, vue graphique, propriétés et onglets bas}{Organisation de l'interface}

\section{Arbre du modèle}

L'arbre du modèle regroupe les objets du projet :

\begin{itemize}
  \item nœuds ;
  \item éléments de barre ;
  \item surfaces ;
  \item matériaux ;
  \item sections ;
  \item cas de charge ;
  \item combinaisons.
\end{itemize}

La sélection d'un objet dans l'arbre synchronise le panneau de propriétés et la
vue graphique.

\section{Vue graphique}

La vue graphique affiche la géométrie du modèle et permet de sélectionner les
objets. Selon les options actives, elle peut afficher :

\begin{itemize}
  \item les nœuds ;
  \item les barres ;
  \item les surfaces ;
  \item les appuis ;
  \item les numéros de nœuds ;
  \item les noms de sections ;
  \item les sections extrudées.
\end{itemize}

\section{Panneau de propriétés}

Le panneau de propriétés affiche les informations de l'objet sélectionné. Il
permet notamment de modifier certaines valeurs directement, par exemple les
coordonnées d'un nœud ou la section affectée à une barre.

\section{Onglets inférieurs}

Les onglets inférieurs contiennent :

\begin{itemize}
  \item la console de messages ;
  \item les résultats ;
  \item le tableau des nœuds ;
  \item le tableau des combinaisons.
\end{itemize}

\section{Menus principaux}

\begin{longtable}{p{0.22\textwidth}p{0.68\textwidth}}
  \toprule
  Menu & Rôle \\
  \midrule
  \menu{Fichier} & Créer, ouvrir, enregistrer un projet. \\
  \menu{Édition} & Annuler, rétablir, copier une sélection. \\
  \menu{Modèle} & Ajouter nœuds, barres, surfaces, matériaux, sections et appuis. \\
  \menu{Charges} & Gérer les cas de charge, charges affectées et combinaisons. \\
  \menu{Analyse} & Choisir le moteur et lancer le calcul. \\
  \menu{Résultats} & Afficher tableaux, diagrammes, déformée et cartes plaques. \\
  \menu{Vue} & Configurer l'affichage et les options graphiques. \\
  \bottomrule
  \caption{Menus principaux}
\end{longtable}
